% Generated draft: requires mathematical and source review.
\documentclass[11pt]{article}
\usepackage[margin=25mm]{geometry}
\usepackage{fontspec}
\setmainfont{Latin Modern Roman}
\usepackage{amsmath,amssymb,mathtools,unicode-math}
\setmathfont{Latin Modern Math}
\usepackage{enumitem,needspace,xurl,hyperref}
\setlength{\parindent}{0pt}
\setlength{\parskip}{5pt}
\setlength{\emergencystretch}{3em}
\newcommand{\N}{\mathbb{N}}
\newcommand{\Z}{\mathbb{Z}}
\newcommand{\Q}{\mathbb{Q}}
\newcommand{\R}{\mathbb{R}}
\newcommand{\C}{\mathbb{C}}
\newcommand{\F}{\mathbb{F}}
\newcommand{\A}{\mathbb{A}}
\newcommand{\PP}{\mathbb{P}}
\newcommand{\E}{\mathbb{E}}
\newcommand{\eps}{\varepsilon}
\newcommand{\dd}{\,\mathrm{d}}
\begin{document}
\section*{UnsolvedMath ID: 30000004}
\textbf{Source problem code:} catalogue entry
\subsection{Definitions and mathematical statement}
Let $I=[0,1]$ and let $M$ be a smooth three-dimensional manifold with nonempty boundary. Fix the endpoint and tangent data in the source's long-knot problem. Write $\operatorname{Emb}(I,M;\partial)$ for the space of smooth embeddings satisfying those prescribed boundary data, with its smooth topology. Let $X^\bullet$ denote the compactified-configuration cosimplicial space associated with this embedding-calculus construction and the same boundary data. Its homotopy totalization is
\[
 \operatorname{Tot}(X^\bullet)=\operatorname{holim}_{[r]\in\Delta}X^r,
\]
interpreted through a suitable fibrant model; $\Delta$ is the simplex category.

The question is whether, for every such $M$ and prescribed admissible endpoint tangent data, the natural embedding-calculus map
\[
 \operatorname{Emb}(I,M;\partial)\longrightarrow\operatorname{Tot}(X^\bullet)
\]
is a homotopy equivalence in dimension three. Here homotopy equivalence means existence of a continuous inverse up to homotopy, as requested by the verified clean statement. The extract does not specify the precise compactification, cosimplicial structure maps, or whether the boundary convention fixes only endpoint jets or whole collar germs; these choices must be matched in the source rather than replaced by a newly chosen configuration model. No compactness, simple-connectivity, or additional dimension hypothesis on $M$ is supplied.

The literature assessment checked on 21 August 2026 reports important dimension-three progress, including surjectivity results for embedding-calculus knot invariants in three-manifolds. It explicitly says the extracted theorem is truncated and does not identify whether the full generalization above is proved. Surjectivity of invariants alone does not establish the asserted homotopy equivalence. The partially-solved label is dated curation, requiring independent verification.
\subsection{Short English statement}
Does the compactified-configuration embedding-calculus model recover the entire homotopy type of long knots in every three-manifold with boundary? The supplied assessment reports partial dimension-three progress but does not identify a proof of this full equivalence.
\subsection{Sources}
\begin{itemize}
\item \emph{The Goodwillie Calculus of Functors} (2005), Oberwolfach Reports, DOI 10.4171/owr/2004/17, pp. 873--914, \url{https://doi.org/10.4171/owr/2004/17}. These title/year and DOI details are reproduced as supplied; the record's proposed year is 2004.
\item Literature sources supplied in the assessment checked 21 August 2026: ``embedding calculus and grope cobordism,'' \url{https://arxiv.org/abs/2010.05120}; ``geometric approach to embedding-calculus invariants,'' \url{https://pure.mpg.de/pubman/item/item_3265796}.
\item Local UnsolvedMath registry, permanent ID $30000004$, exact source code \texttt{catalogue entry}, created 25 August 2026 and updated 16 September 2026. Its clean statement is labelled corrected and verified against the official report; its literature triage is dated and does not establish the full claimed equivalence.
\end{itemize}
\end{document}
